% GENERATED FILE. Edit canonical lessons, then run npm run generate:books.
% canonical-source-hash: fnv1a64:1e4d9d32b02c99e3
% canonical-lessons: JA-C464-kankyaku, JA-C464-shinpan, JA-C464-tokuten, JA-C464-semeru, JA-C464-ubau

\chapter{At the Match}
\label{ch:ja-at-the-match}
\hlchaptermodality{464}

\begin{chapteropening}
\textbf{By the end of this chapter:} \emph{I can talk about a match I went to.}

Complete the last lesson of chapter 464: I can talk about a match I went to.
\end{chapteropening}

\section[\texorpdfstring{\ja{かんきゃく} (kankyaku)}{かんきゃく (kankyaku)}]{\ja{かんきゃく} (kankyaku) — spectators; the audience}

\label{lesson:JA-C464-kankyaku}



\begin{quote}
Before the new one: say the Japanese for to bounce, then the Japanese for to vie.
\end{quote}

\subsection*{You'll want to know: \ja{かんきゃく}}
\textbf{\ja{かんきゃく}} — \emph{kankyaku} — \textquotedblleft{}spectators; the audience\textquotedblright{}.

\begin{quote}
\textbf{\ja{かんきゃくが} \ja{おうえんしています。}} — \emph{kankyaku ga ōen shite imasu} — \textquotedblleft{}The spectators are cheering.\textquotedblright{}
\end{quote}

\begin{cousinweb}
\textbf{\ja{かんきゃく}} (\emph{kankyaku}) is Sino-Japanese (kango): two Chinese characters read with their on'yomi, readings borrowed from Chinese. They are \emph{kan}, \textquotedblleft{}view\textquotedblright{} (Middle Chinese \emph{kwan}) and \emph{kyaku}, \textquotedblleft{}guest\textquotedblright{} (\emph{khaek}) (KANJIDIC2; Middle Chinese in Baxter's transcription).
\end{cousinweb}

\begin{grammarlens}[title={What you've built}]
One more word for describing what you lived and felt.
\end{grammarlens}

\subsection*{Your turn}
\begin{itemize}
\raggedright
  \item \emph{Say it:} \emph{kankyaku}
  \item \emph{Say it:} \emph{kankyaku}, once more
  \item \emph{Say it:} \emph{kankyaku ga ōen shite imasu} — The spectators are cheering.
  \item \emph{Recall:} say the Japanese for to bounce, then the Japanese for to vie, then say \emph{kankyaku} again
\end{itemize}

\begin{tcolorbox}[breakable,colback=teal!4,colframe=teal!35!black,title={Before you move on}]
What does \ja{かんきゃく} mean? (\textquotedblleft{}Spectators; the audience\textquotedblright{}.) Say it once more.
\end{tcolorbox}
\section[\texorpdfstring{\ja{しんぱん} (shinpan)}{しんぱん (shinpan)}]{\ja{しんぱん} (shinpan) — a referee; an umpire}

\label{lesson:JA-C464-shinpan}



\begin{quote}
Before the new one: say the Japanese for to vie, then the Japanese for spectators.
\end{quote}

\subsection*{You'll want to know: \ja{しんぱん}}
\textbf{\ja{しんぱん}} — \emph{shinpan} — \textquotedblleft{}a referee; an umpire\textquotedblright{}.

\begin{quote}
\textbf{\ja{しんぱんが} \ja{ふえを} \ja{ふきました。}} — \emph{shinpan ga fue o fukimashita} — \textquotedblleft{}The referee blew the whistle.\textquotedblright{}
\end{quote}

\begin{cousinweb}
\textbf{\ja{しんぱん}} (\emph{shinpan}) is Sino-Japanese (kango): two Chinese characters read with their on'yomi, readings borrowed from Chinese. They are \emph{shin}, \textquotedblleft{}judge\textquotedblright{} (Middle Chinese \emph{syimX}) and \emph{pan}, \textquotedblleft{}judgement\textquotedblright{} (\emph{phanH}) (KANJIDIC2; Middle Chinese in Baxter's transcription).
\end{cousinweb}

\begin{grammarlens}[title={What you've built}]
One more word for describing what you lived and felt.
\end{grammarlens}

\subsection*{Your turn}
\begin{itemize}
\raggedright
  \item \emph{Say it:} \emph{shinpan}
  \item \emph{Say it:} \emph{shinpan}, once more
  \item \emph{Say it:} \emph{shinpan ga fue o fukimashita} — The referee blew the whistle.
  \item \emph{Recall:} say the Japanese for to vie, then the Japanese for spectators, then say \emph{shinpan} again
\end{itemize}

\begin{tcolorbox}[breakable,colback=teal!4,colframe=teal!35!black,title={Before you move on}]
What does \ja{しんぱん} mean? (\textquotedblleft{}A referee; an umpire\textquotedblright{}.) Say it once more.
\end{tcolorbox}
\section[\texorpdfstring{\ja{とくてん} (tokuten)}{とくてん (tokuten)}]{\ja{とくてん} (tokuten) — points scored; a score}

\label{lesson:JA-C464-tokuten}



\begin{quote}
Before the new one: say the Japanese for spectators, then the Japanese for a referee.
\end{quote}

\subsection*{You'll want to know: \ja{とくてん}}
\textbf{\ja{とくてん}} — \emph{tokuten} — \textquotedblleft{}points scored; a score\textquotedblright{}.

\begin{quote}
\textbf{\ja{せんしゅが} \ja{さいごに} \ja{とくてんしました。}} — \emph{senshu ga saigo ni tokuten shimashita} — \textquotedblleft{}The player scored at the end.\textquotedblright{}
\end{quote}

\begin{cousinweb}
\textbf{\ja{とくてん}} (\emph{tokuten}) is Sino-Japanese (kango): two Chinese characters read with their on'yomi, readings borrowed from Chinese. They are \emph{toku}, \textquotedblleft{}gain\textquotedblright{} (Middle Chinese \emph{tok}) and \emph{ten}, \textquotedblleft{}point\textquotedblright{} (\emph{temX}) (KANJIDIC2; Middle Chinese in Baxter's transcription).
\end{cousinweb}

\begin{grammarlens}[title={What you've built}]
One more word for describing what you lived and felt.
\end{grammarlens}

\subsection*{Your turn}
\begin{itemize}
\raggedright
  \item \emph{Say it:} \emph{tokuten}
  \item \emph{Say it:} \emph{tokuten}, once more
  \item \emph{Say it:} \emph{senshu ga saigo ni tokuten shimashita} — The player scored at the end.
  \item \emph{Recall:} say the Japanese for spectators, then the Japanese for a referee, then say \emph{tokuten} again
\end{itemize}

\begin{tcolorbox}[breakable,colback=teal!4,colframe=teal!35!black,title={Before you move on}]
What does \ja{とくてん} mean? (\textquotedblleft{}Points scored; a score\textquotedblright{}.) Say it once more.
\end{tcolorbox}
\section[\texorpdfstring{\ja{せめる} (semeru)}{せめる (semeru)}]{\ja{せめる} (semeru) — to attack (in a game)}

\label{lesson:JA-C464-semeru}



\begin{quote}
Before the new one: say the Japanese for a referee, then the Japanese for points scored.
\end{quote}

\subsection*{You'll want to know: \ja{せめる}}
\textbf{\ja{せめる}} — \emph{semeru} — \textquotedblleft{}to attack (in a game)\textquotedblright{}.

\begin{quote}
\textbf{\ja{さいごまで} \ja{せめて} \ja{ください。}} — \emph{saigo made semete kudasai} — \textquotedblleft{}Please keep attacking to the end.\textquotedblright{}
\end{quote}

Polite \textbf{\ja{せめます}}, \emph{sememasu}; an ichidan verb. Past (polite): \textbf{\ja{せんしゅが} \ja{せめました。}}, \emph{senshu ga sememashita}.

\begin{cousinweb}
\textbf{\ja{せめる}} (\emph{semeru}) is a native Japanese word (wago), inherited rather than borrowed from Chinese; its kanji spelling was attached to it later (JMdict).
\end{cousinweb}

\begin{grammarlens}[title={What you've built}]
One more word for describing what you lived and felt.
\end{grammarlens}

\subsection*{Your turn}
\begin{itemize}
\raggedright
  \item \emph{Say it:} \emph{semeru}
  \item \emph{Say it:} \emph{semeru}, once more
  \item \emph{Say it:} \emph{saigo made semete kudasai} — Please keep attacking to the end.
  \item \emph{Recall:} say the Japanese for a referee, then the Japanese for points scored, then say \emph{semeru} again
\end{itemize}

\begin{tcolorbox}[breakable,colback=teal!4,colframe=teal!35!black,title={Before you move on}]
What does \ja{せめる} mean? (\textquotedblleft{}To attack (in a game)\textquotedblright{}.) Say it once more.
\end{tcolorbox}
\section[\texorpdfstring{\ja{うばう} (ubau)}{うばう (ubau)}]{\ja{うばう} (ubau) — to snatch; to take away}

\label{lesson:JA-C464-ubau}



\begin{quote}
Before the new one: say the Japanese for points scored, then the Japanese for to attack.
\end{quote}

\subsection*{You'll want to know: \ja{うばう}}
\textbf{\ja{うばう}} — \emph{ubau} — \textquotedblleft{}to snatch; to take away\textquotedblright{}.

\begin{quote}
\textbf{\ja{あめが} \ja{かんきゃくの} \ja{たのしみを} \ja{うばいました。}} — \emph{ame ga kankyaku no tanoshimi o ubaimashita} — \textquotedblleft{}The rain took away the spectators' fun.\textquotedblright{}
\end{quote}

Polite \textbf{\ja{うばいます}}, \emph{ubaimasu}; a godan verb. Past (polite): \textbf{\ja{せんしゅが} \ja{まりを} \ja{うばいました。}}, \emph{senshu ga mari o ubaimashita}.

\begin{cousinweb}
\textbf{\ja{うばう}} (\emph{ubau}) is a native Japanese word (wago), inherited rather than borrowed from Chinese; its kanji spelling was attached to it later (JMdict).
\end{cousinweb}

\begin{grammarlens}[title={What you've built}]
One more word for describing what you lived and felt.
\end{grammarlens}

\subsection*{Your turn}
\begin{itemize}
\raggedright
  \item \emph{Say it:} \emph{ubau}
  \item \emph{Say it:} \emph{ubau}, once more
  \item \emph{Say it:} \emph{ame ga kankyaku no tanoshimi o ubaimashita} — The rain took away the spectators' fun.
  \item \emph{Recall:} say the Japanese for points scored, then the Japanese for to attack, then say \emph{ubau} again
\end{itemize}

\begin{tcolorbox}[breakable,colback=teal!4,colframe=teal!35!black,title={Before you move on}]
What does \ja{うばう} mean? (\textquotedblleft{}To snatch; to take away\textquotedblright{}.) Say it once more.
\end{tcolorbox}
